\roteiro{1}

% Slide 3: persistência poliglota (relatório §1.1, §2.1, Fig. 1)
% A Fig. 1 do relatório (images/figure1-polyglot-ecommerce.png) rende texto de ~4,7 pt em
% largura total; o desenho abaixo repete o seu conteúdo em TikZ, com texto legível.
\begin{frame}{Persistência poliglota}
  \begin{itemize}
    \item Um só modelo de dados não atende todas as necessidades.
    \item \destaque{Persistência poliglota}: cada tipo de dado fica no SGBD (sistema
          gerenciador de banco de dados) mais adequado.
  \end{itemize}
  \vfill
  \centering
  \begin{tikzpicture}[x=1mm, y=1mm, font=\footnotesize,
      caixa/.style={draw=ufscborda, fill=white, rounded corners=0.6mm, inner sep=2pt,
                    minimum height=6.2mm, anchor=west, align=left, text=ufsctexto},
      seta/.style={->, >=stealth, draw=ufscazul!70, thick}]
    \node[caixa, fill=ufscfundo, align=center, minimum width=34mm] (csv) at (0,0)
      {CSV operacional\\[-1pt]{\scriptsize\texttt{ecommerce\_join.csv}}};
    \node[caixa, fill=ufscazul, draw=ufscazul, text=white, align=center, minimum width=34mm]
      (tool) at (42,0) {\bfseries Polyglot Import CSV\\[-1pt]{\scriptsize configuração JSON}};
    \node[caixa, minimum width=58mm] (n1) at (88, 14)  {\textbf{PostgreSQL} · pedidos e inventário};
    \node[caixa, minimum width=58mm] (n2) at (88, 7)   {\textbf{MongoDB} · catálogo de produtos};
    \node[caixa, minimum width=58mm] (n3) at (88, 0)   {\textbf{Redis} · carrinho e sessão};
    \node[caixa, minimum width=58mm] (n4) at (88, -7)  {\textbf{Cassandra} · log de atividade};
    \node[caixa, minimum width=58mm] (n5) at (88, -14) {\textbf{Neo4j} · recomendações};
    \draw[seta] (csv.east) -- (tool.west);
    \foreach \i in {1,...,5} \draw[seta] (tool.east) -- (n\i.west);
  \end{tikzpicture}\par
  \medskip
  {\footnotesize Cada dado do e-commerce no SGBD mais adequado.
   CSV: texto com valores separados por vírgula.\par}
\end{frame}

% Slide 4: o problema (relatório §1.1, §3.1)
\begin{frame}{O problema}
  \begin{itemize}
    \item Cada SGBD tem o seu \destaque{importador nativo}, com sintaxe própria.
    \item Cada importador atende um único SGBD por vez.
    \item Nenhuma ferramenta encontrada importa o mesmo CSV nos cinco.
  \end{itemize}
  \vfill
  {\footnotesize\bfseries\color{ufscazul} Importadores nativos de CSV\par}
  \smallskip
  \footnotesize
  \setlength{\tabcolsep}{3pt}
  \renewcommand{\arraystretch}{1.2}
  \begin{tabular}{@{}E{0.1\textwidth}*{5}{E{0.16\textwidth}}@{}}
    \toprule
    & \bfseries PostgreSQL & \bfseries Redis & \bfseries MongoDB & \bfseries Cassandra
    & \bfseries Neo4j \\
    \midrule
    \textcolor{black!60}{Comando}
      & \texttt{COPY}
      & \texttt{redis-cli -{}-pipe}
      & \texttt{mongoimport}
      & \texttt{cqlsh} e \texttt{sstableloader}
      & \texttt{neo4j-admin import} \\
    \textcolor{black!60}{Detalhe}
      & colunas na ordem do CSV
      & arquivo em protocolo bruto
      & tipo e cabeçalho por parâmetros
      & conversão para SSTable (formato próprio)
      & CSV separado para nós e relacionamentos \\
    \bottomrule
  \end{tabular}
\end{frame}

% Slide 5: objetivos (relatório §1.2)
\begin{frame}{Objetivos}
  \begin{painel}{Objetivo geral}{}
    Desenvolver a ferramenta \textit{Polyglot Import CSV}, que importa CSV para
    PostgreSQL, Redis, MongoDB, Cassandra ou Neo4j.
  \end{painel}
  \smallskip
  \begin{painel}{Objetivos específicos}{}
    \small
    \begin{enumerate}
      \item \destaque{Configuração} em JSON (texto estruturado): filtros, chave primária e
            nome da coluna.
      \item \destaque{Validação prévia}: um erro na configuração impede o início da
            importação.
      \item \destaque{Avaliação} do desempenho com diferentes volumes e modelos de dados.
      \item \destaque{Acesso} por linha de comando (CLI) e por interface gráfica (GUI).
    \end{enumerate}
  \end{painel}
\end{frame}
